\documentclass[11pt]{article}
\usepackage[margin=1in]{geometry}
\usepackage{amsmath,amssymb,amsthm}
\usepackage{graphicx}
\usepackage[hidelinks]{hyperref}

\newtheorem{theorem}{Theorem}
\newtheorem{corollary}{Corollary}
\newtheorem{remark}{Remark}

\newcommand{\R}{\mathbb R}
\newcommand{\Prob}{\mathrm{Prob}}
\newcommand{\OO}{\mathcal O}

\title{A Product Counterexample to Arbitrarily Slow Proper \(L^p\)-Cocycles}
\author{fa\_banach\_001 candidate packet}
\date{2026-07-03}

\begin{document}
\maketitle

\section*{Status}

Status: \textbf{counterexample\_likely\_valid}.

This packet concerns Question 1 in Section 6 of Antonio Lopez Neumann and Juan
Paucar, \emph{On growth of cocycles of isometric representations on
\(L^p\)-spaces}, arXiv:2501.12808v2.

\section*{Source Question}

The source paper defines a compactly generated locally compact group \(G\) to
have arbitrarily slow proper cocycles in \(L^p\) if, for every proper function
\(f:G\to [0,\infty)\), there are an \(L^p\)-space \(E\), a continuous
isometric representation \(\pi:G\to \OO(E)\), and a proper cocycle
\(b\in Z^1(G,\pi)\) such that \(\|b(g)\|\leq f(g)\) for every \(g\in G\).

Question 1 asks:

\begin{quote}
Let \(G\) be a locally compact compactly generated group. Suppose that \(G\)
is \(a\)-\(FL^p\)-menable (i.e. \(G\) admits a proper cocycle on an
\(L^p\)-space) for some \(2<p<\infty\), not \(a\)-\(T\)-menable and does not
have property \((T)\). Does \(G\) have arbitrarily slow proper cocycles in
\(L^p\)?
\end{quote}

The question appears on page 28 of the locally rendered source PDF:

\begin{center}
\includegraphics[width=\textwidth,height=0.82\textheight,keepaspectratio]{figures/open_problem_crop.png}
\end{center}

\section*{Proof Intuition}

The obstruction is already present in the product examples mentioned by the
source paper. Put a property \((T)\) factor \(B\) next to a non-\((T)\)
\(a\)-\(FL^p\)-menable factor \(A\). The product \(A\times B\) is still
\(a\)-\(FL^p\)-menable and is not \(a\)-\(T\)-menable, but any proper cocycle
on the product must be unbounded on the \(B\)-axis. Now choose the demanded
``arbitrarily slow'' majorant to grow only logarithmically in that \(B\)
direction. The source paper's property-\((T)\) theorem forces every unbounded
\(L^p\)-cocycle of \(B\) to have random-walk \(L^p\)-average at least
\(n^{1/p}\), contradicting the logarithmic upper bound.

\section*{Counterexample}

\begin{theorem}
Fix \(2<p<\infty\). Let \(B\) be a noncompact compactly generated locally
compact group with property \((T)\) that is \(a\)-\(FL^p\)-menable. Let \(A\)
be a noncompact compactly generated locally compact group without property
\((T)\) that is \(a\)-\(FL^p\)-menable. Then \(G=A\times B\) satisfies the
hypotheses of Question 1 for this \(p\), but \(G\) does not have arbitrarily
slow proper cocycles in \(L^p\).
\end{theorem}

\begin{proof}
The product \(G=A\times B\) is compactly generated. Since both factors admit
proper \(L^p\)-cocycles, their \(\ell_p\)-direct sum is a proper
\(L^p\)-cocycle on \(G\), so \(G\) is \(a\)-\(FL^p\)-menable. The quotient
map \(G\to A\) shows that \(G\) does not have property \((T)\). Also \(G\) is
not \(a\)-\(T\)-menable: otherwise the closed subgroup \(B\) would be
\(a\)-\(T\)-menable; combined with property \((T)\), this would force \(B\)
to be compact, contradicting the hypothesis.

Choose compact generating sets \(S_A,S_B\) and let \(|\cdot|_A,|\cdot|_B\)
be the associated word lengths. Define
\[
  f(a,b)=1+|a|_A+\log(1+|b|_B),\qquad (a,b)\in A\times B .
\]
This is a proper function on \(G\), because its sublevel sets are contained in
products of word balls in \(A\) and \(B\).

Assume, toward a contradiction, that \(G\) has arbitrarily slow proper
\(L^p\)-cocycles. Then there are an \(L^p\)-space \(E\), a continuous
isometric representation \(\pi:G\to \OO(E)\), and a proper cocycle
\(c\in Z^1(G,\pi)\) such that
\[
  \|c(a,b)\|\leq f(a,b)\qquad\text{for all }(a,b)\in G .
\]
Restrict to the \(B\)-factor:
\[
  c_B(b)=c(e_A,b),\qquad b\in B .
\]
This is a \(1\)-cocycle for the restricted representation of \(B\). It is
proper, hence unbounded, because the inclusion \(B\to A\times B\),
\(b\mapsto(e_A,b)\), is proper and \(c\) is proper.
In particular \(c_B\) is not a coboundary, since every coboundary
\(b\mapsto v-\pi(b)v\) is bounded.

By the separable reduction in Appendix A of the source paper, the restricted
cocycle may be regarded inside a separable \(L^p\)-subspace invariant under
\(B\). Since \(B\) has property \((T)\) and is \(a\)-\(FL^p\)-menable, it
does not have property \(FL^p\). The source paper's Theorem D therefore
applies to the unbounded cocycle \(c_B\): for every compactly supported
cohomologically adapted probability measure \(\mu\) on \(B\), there is a
constant \(C>0\) such that
\[
  \|c_B\|_{L^p(\mu^{*n})}\geq C^{-1}n^{1/p}
  \qquad\text{for all }n\geq 1 .
\]

Choose such a measure \(\mu\) supported in a compact generating set contained
in \(S_B\). Then \(\mu^{*n}\) is supported in the word ball
\(\{b\in B: |b|_B\leq n\}\). The assumed pointwise bound gives
\[
  \|c_B(b)\|\leq 1+\log(1+|b|_B)\leq 1+\log(1+n)
  \qquad (b\in \operatorname{supp}\mu^{*n}),
\]
and hence
\[
  \|c_B\|_{L^p(\mu^{*n})}\leq 1+\log(1+n).
\]
This contradicts \(C^{-1}n^{1/p}\leq 1+\log(1+n)\) for large \(n\). Thus
the proper function \(f\) is not dominated by any proper \(L^p\)-cocycle on
\(G\), and \(G\) does not have arbitrarily slow proper cocycles in \(L^p\).
\end{proof}

\begin{corollary}
For suitable sufficiently large \(p\), the product groups
\(\mathrm{SO}(n,1)\times \mathrm{Sp}(m,1)\), \(n,m\geq 2\), are negative
answers to Question 1.
\end{corollary}

\begin{proof}
The source paper itself lists these products as examples satisfying the
hypotheses of Question 1. More explicitly, \(\mathrm{SO}(n,1)\) is
\(a\)-\(T\)-menable and does not have property \((T)\), while
\(\mathrm{Sp}(m,1)\) is noncompact and has property \((T)\). For sufficiently
large \(p\), the rank-one hyperbolic factor \(\mathrm{Sp}(m,1)\) admits a
proper affine isometric action on an \(L^p\)-space by the standard
Bourdon--Pajot/Pansu \(L^p\)-cohomology construction cited in the source
paper. The theorem applies.
\end{proof}

\section*{Verification Notes}

The argument uses no computation. The only growth comparison is the elementary
asymptotic contradiction \(n^{1/p}\not=O(\log n)\).

Duplicate and novelty checks were bounded as follows. I searched the cheap run
indexes for arXiv:2501.12808, ``arbitrarily slow proper cocycles'',
``slow proper cocycles'', ``a-FL\^p'', and the product examples; no prior
packet or exact duplicate was found. Web searches on 2026-07-03 for the exact
question and close variants found the source paper, arXiv:2501.12808, but no
later exact resolution. The novelty confidence is therefore
\emph{likely valid / needs human review}, not human-verified.

\section*{Human Review Recommendation}

Check two points first: (1) that the source paper's Theorem D is being applied
with the intended separable reduction to the restricted \(B\)-cocycle, and
(2) that the definition of ``arbitrarily slow proper cocycles in \(L^p\)''
allows arbitrary proper control functions such as
\(1+|a|_A+\log(1+|b|_B)\). If both are accepted, the contradiction is direct.

\begin{thebibliography}{9}

\bibitem{LNP}
Antonio Lopez Neumann and Juan Paucar.
\newblock \emph{On growth of cocycles of isometric representations on
\(L^p\)-spaces}.
\newblock arXiv:2501.12808v2, 2026.

\bibitem{BFGM}
Uri Bader, Alex Furman, Tsachik Gelander, and Nicolas Monod.
\newblock Property \((T)\) and rigidity for actions on Banach spaces.
\newblock \emph{Acta Mathematica} 198 (2007), 57--105.

\bibitem{BP}
Marc Bourdon and Herve Pajot.
\newblock Cohomologie \(l_p\) et espaces de Besov.
\newblock \emph{Journal fur die reine und angewandte Mathematik} 558
(2003), 85--108.

\bibitem{BHV}
Bachir Bekka, Pierre de la Harpe, and Alain Valette.
\newblock \emph{Kazhdan's Property (T)}.
\newblock Cambridge University Press, 2008.

\end{thebibliography}

\end{document}
